\chapter{2006年山东卷（理）}

\section{单选题}

\begin{problem}
定义集合运算：\(A\odot B=\{z\mid z=xy(x+y),\ x\in A,\ y\in B\}\)，设集合\(A=\{0,1\}\)，\(B=\{2,3\}\)，则集合\(A\odot B\)的所有元素之和为（\quad）

\choices
    {\(0\)}
    {\(6\)}
    {\(12\)}
    {\(18\)}
\end{problem}

\begin{answer}
D
\end{answer}

\begin{solution}
当\(x=0\)时，\(z=0\)；当\(x=1\)时，分别取\(y=2\)，\(3\)，得到\(z=1\cdot2(1+2)=6\)和\(z=1\cdot3(1+3)=12\). 因此\(A\odot B=\{0,6,12\}\)，所有元素之和为\(0+6+12=18\)，故选 D.
\end{solution}

\begin{problem}
函数\(y=1+a^x\)（\(0<a<1\)）的反函数的图象大致是（\quad）

\choices
    {\choicebitmap[width=0.15\paperwidth]{img/2006/shandong_science/q2_choice_a.png}}
    {\choicebitmap[width=0.15\paperwidth]{img/2006/shandong_science/q2_choice_b.png}}
    {\choicebitmap[width=0.15\paperwidth]{img/2006/shandong_science/q2_choice_c.png}}
    {\choicebitmap[width=0.15\paperwidth]{img/2006/shandong_science/q2_choice_d.png}}
\end{problem}

\begin{answer}
A
\end{answer}

\begin{solution}
由\(0<a<1\)，函数\(y=1+a^x\)单调递减，值域为\((1,+\infty)\)，且图象经过点\((0,2)\). 反函数的图象是原函数图象关于直线\(y=x\)的对称图象，所以反函数的定义域为\((1,+\infty)\)，并经过点\((2,0)\)，同时仍然单调递减. 当\(x\to1^+\)时，反函数值趋于\(+\infty\)，故其图象具有竖直渐近线\(x=1\)，并经过\((2,0)\)，符合选项 A.
\end{solution}

\begin{problem}
设\(f(x)=\begin{cases}
2\e^{x-1},&x<2,\\
\log_3(x^2-1),&x\geqslant2,
\end{cases}\) 则不等式\(f(x)>2\)的解集为（\quad）

\choices
    {\((1,2)\cup(3,+\infty)\)}
    {\((\sqrt{10},+\infty)\)}
    {\((1,2)\cup(\sqrt{10},+\infty)\)}
    {\((1,2)\)}
\end{problem}

\begin{answer}
C
\end{answer}

\begin{solution}
当\(x<2\)时，\(2\e^{x-1}>2\)等价于\(\e^{x-1}>1\)，即\(x>1\)，所以这一段的解为\((1,2)\). 当\(x\geqslant2\)时，\(\log_3(x^2-1)>2\)等价于\(x^2-1>3^2\)，即\(x^2>10\). 结合\(x\geqslant2\)，得到\(x>\sqrt{10}\). 因此原不等式的解集为\((1,2)\cup(\sqrt{10},+\infty)\)，故选 C.
\end{solution}

\begin{problem}
在\(\triangle ABC\)中，角\(A\)，\(B\)，\(C\)的对边分别为\(a\)，\(b\)，\(c\)，已知\(A=\dfrac{\pi}{3}\)，\(a=\sqrt{3}\)，\(b=1\)，则\(c=\)（\quad）

\choices
    {\(1\)}
    {\(2\)}
    {\(\sqrt{3}-1\)}
    {\(\sqrt{3}\)}
\end{problem}

\begin{answer}
B
\end{answer}

\begin{solution}
由余弦定理，\(a^2=b^2+c^2-2bc\cos A\). 代入\(a=\sqrt{3}\)，\(b=1\)，\(A=\dfrac{\pi}{3}\)，得\(3=1+c^2-c\)，
即\(c^2-c-2=0\)，所以\(c=2\)或\(c=-1\). 因为边长\(c>0\)，故\(c=2\)，选 B.
\end{solution}

\begin{problem}
设向量\(\bs{a}=(1,-2)\)，\(\bs{b}=(-2,4)\)，\(\bs{c}=(-1,-2)\). 若表示向量\(4\bs{a}\)，\(4\bs{b}-2\bs{c}\)，\(2(\bs{a}-\bs{c})\)，\(\bs{d}\)的有向线段首尾相接能构成四边形，则向量\(\bs{d}\)为（\quad）

\choices
    {\((2,6)\)}
    {\((-2,6)\)}
    {\((2,-6)\)}
    {\((-2,-6)\)}
\end{problem}

\begin{solution}
四条有向线段首尾相接构成封闭四边形时，沿四条有向线段走一周，位移和为零向量，因此
\[
\bs{d}=-4\bs{a}-(4\bs{b}-2\bs{c})-2(\bs{a}-\bs{c})
      =-6\bs{a}-4\bs{b}+4\bs{c}
\]
由题面给出的坐标，
\(-6\bs{a}=(-6,12)\)，\(-4\bs{b}=(8,-16)\)，\(4\bs{c}=(-4,-8)\)，从而\(\bs{d}=(-6,12)+(8,-16)+(-4,-8)=(-2,-12)\).
该结果不在原始试卷题面给出的四个选项中，所以按原题面不能确定一个合法选项，现保留这一客观题面冲突，不臆造选择项.
\end{solution}

\begin{problem}
已知定义在\(\R\)上的奇函数\(f(x)\)满足\(f(x+2)=-f(x)\)，则\(f(6)\)的值为（\quad）

\choices
    {\(-1\)}
    {\(0\)}
    {\(1\)}
    {\(2\)}
\end{problem}

\begin{answer}
B
\end{answer}

\begin{solution}
由\(f(x+2)=-f(x)\)，依次取\(x=4\)，\(2\)，\(0\)，得
\[
f(6)=-f(4)=f(2)=-f(0)
\]
奇函数满足\(f(0)=0\)，所以\(f(6)=0\)，故选 B.
\end{solution}

\begin{problem}
在给定椭圆中，过焦点且垂直于长轴的弦长为\(\sqrt{2}\)，焦点到相应准线的距离为\(1\)，则该椭圆的离心率为（\quad）

\choices
    {\(\sqrt{2}\)}
    {\(\dfrac{\sqrt{2}}{2}\)}
    {\(\dfrac{1}{2}\)}
    {\(\dfrac{\sqrt{3}}{2}\)}
\end{problem}

\begin{answer}
B
\end{answer}

\begin{solution}
设椭圆的长半轴、短半轴和半焦距分别为\(a\)，\(b\)，\(c\)，则\(c^2=a^2-b^2\). 过焦点且垂直于长轴的弦为通径，所以
\(\dfrac{2b^2}{a}=\sqrt{2}\)，且\(\dfrac{b^2}{c}=1\).
由第二式得\(b^2=c\)，再由第一式得\(c=\dfrac{a}{\sqrt{2}}\). 代入\(c^2=a^2-b^2\)，得到
\(\dfrac{a^2}{2}=a^2-\dfrac{a}{\sqrt{2}}\)，
所以\(a=\sqrt{2}\)，\(c=1\). 因而离心率
\[
e=\frac{c}{a}=\frac{1}{\sqrt{2}}=\frac{\sqrt{2}}{2}
\]
故选 B.
\end{solution}

\begin{problem}
设\(p:x^2-x-20>0\)，\(q:\dfrac{1-x^2}{|x|-2}<0\)，则\(p\)是\(q\)的（\quad）

\choices
    {充分不必要条件}
    {必要不充分条件}
    {充要条件}
    {既不充分也不必要条件}
\end{problem}

\begin{answer}
A
\end{answer}

\begin{solution}
因为\(x^2-x-20=(x-5)(x+4)\)，
所以\(p\)的解集为\((-\infty,-4)\cup(5,+\infty)\). 令\(t=|x|\)，则\(\dfrac{1-x^2}{|x|-2}=\dfrac{1-t^2}{t-2}\).
当\(0\leqslant t<1\)时，分子为正、分母为负，分式小于\(0\)；当\(1<t<2\)时，分子、分母均为负，分式大于\(0\)；当\(t>2\)时，分子为负、分母为正，分式小于\(0\). 另外\(t=1\)使分子为\(0\)，\(t=2\)使分母为\(0\)，均不满足不等式. 因此\(q\)的解集为\((-\infty,-2)\cup(-1,1)\cup(2,+\infty)\).
由\((-\infty,-4)\cup(5,+\infty)\subset(-\infty,-2)\cup(2,+\infty)\)，可知\(p\)是\(q\)的充分条件；例如\(x=0\)满足\(q\)而不满足\(p\)，所以\(p\)不是\(q\)的必要条件. 因此\(p\)是\(q\)的充分不必要条件，故选 A.
\end{solution}

\begin{problem}
已知集合\(A=\{5\}\)，\(B=\{1,2\}\)，\(C=\{1,3,4\}\)，从这三个集合各取一个元素构成空间直角坐标系中点的坐标，则确定的不同点的个数为（\quad）

\choices
    {\(33\)}
    {\(34\)}
    {\(35\)}
    {\(36\)}
\end{problem}

\begin{answer}
A
\end{answer}

\begin{solution}
先从\(A\)，\(B\)，\(C\)中各取一个数.若\(B\)取\(1\)，\(C\)取\(1\)，所得数字为\(5\)，\(1\)，\(1\)，将它们安排到三个坐标位置有\(\dfrac{3!}{2!}=3\)种不同点. 其余五种取法所得三个数字互不相同，每种都有\(3!=6\)种安排方法. 因此不同点的个数为
\(3+5\times6=33\)，
故选 A.
\end{solution}

\begin{problem}
已知\(\left(x^{3}-\dfrac{\i}{\sqrt{x}}\right)^n\)的展开式中第三项与第五项的系数之比为\(-\dfrac{3}{14}\)，其中\(\i^2=-1\)，则展开式中常数项是（\quad）

\choices
    {\(-45\i\)}
    {\(45\i\)}
    {\(-45\)}
    {\(45\)}
\end{problem}

\begin{solution}
按原始试卷题面写出的二项式，通项为
\[
T_{r+1}=\mathrm C_n^r(-\i)^r x^{3(n-r)-\frac r2}
       =\mathrm C_n^r(-\i)^r x^{3n-\frac72r}
\]
第三项和第五项分别对应\(r=2\)，\(4\)，其系数之比为
\[
\frac{\mathrm C_n^2(-\i)^2}{\mathrm C_n^4(-\i)^4}
=-\frac{\mathrm C_n^2}{\mathrm C_n^4}
=-\frac{3}{14}
\]
因此
\[
\frac{\mathrm C_n^2}{\mathrm C_n^4}=\frac{3}{14}
\Longrightarrow
\frac{12}{(n-2)(n-3)}=\frac{3}{14}
\Longrightarrow (n-2)(n-3)=56
\]
整理为\(n^2-5n-50=0\)，所以\(n=10\)或\(n=-5\). 由于第三项和第五项存在，\(n\geqslant4\)，故\(n=10\).
此时常数项应满足
\[
3n-\frac72r=0
\Longrightarrow r=\frac{6n}{7}=\frac{60}{7}
\]
但二项式展开中的项指标\(r\)必须是整数，且\(0\leqslant r\leqslant n\). 因此按原始试卷题面二项式展开不存在常数项，题面给出的四个选项均不符合该计算结果，现保留这一客观题面冲突，不臆造答案.
\end{solution}

\begin{problem}
某公司招收男职员\(x\)名，女职员\(y\)名，\(x\)和\(y\)须满足约束条件
\(\begin{cases}
5x-11y\geqslant-22,\\
2x+3y\geqslant9,\\
2x\leqslant11
\end{cases}\)，则\(z=10x+10y\)的最大值是（\quad）

\choices
    {\(80\)}
    {\(85\)}
    {\(90\)}
    {\(95\)}
\end{problem}

\begin{answer}
C
\end{answer}

\begin{solution}
因为\(x\)，\(y\)表示人数，所以为非负整数. 由\(2x\leqslant11\)，得\(x\leqslant5\). 又由第一条约束，
\(11y\leqslant5x+22\leqslant47\)，故\(y\leqslant4\). 因而\(z=10(x+y)\leqslant10(5+4)=90\).
取\(x=5\)，\(y=4\)时，三条约束分别成立，且\(z=90\)，所以最大值为\(90\)，故选 C.
\end{solution}

\begin{problem}
如图，在等腰梯形\(ABCD\)中，\(AB=2DC=2\)，\(\angle DAB=60^\circ\)，\(E\)为\(AB\)的中点，将\(\triangle ADE\)与\(\triangle BEC\)分别沿\(ED\)，\(EC\)向上折起，使\(A\)，\(B\)重合于点\(P\)，则三棱锥\(P-DCE\)的外接球的体积为（\quad）

\bitmapfigure[max width=5.8cm]{img/2006/shandong_science/q12_fig1.png}

\choices
    {\(\dfrac{4\sqrt{3}\pi}{27}\)}
    {\(\dfrac{\sqrt{6}\pi}{2}\)}
    {\(\dfrac{\sqrt{6}\pi}{8}\)}
    {\(\dfrac{\sqrt{6}\pi}{24}\)}
\end{problem}

\begin{answer}
C
\end{answer}

\begin{solution}
由\(AB=2\)，\(DC=1\)，且梯形为等腰梯形，过\(D\)，\(C\)分别向\(AB\)作垂线，则每侧的水平差为\(\dfrac12\). 由\(\angle DAB=60^\circ\)，得
\(AD=BC=\dfrac{\frac12}{\cos60^\circ}=1\).
又\(AE=BE=1\)，由余弦定理
\(DE^2=AD^2+AE^2-2\cdot AD\cdot AE\cos60^\circ=1\)，同理\(CE=1\). 所以折叠前\(\triangle ADE\)和\(\triangle BEC\)都是边长为\(1\)的等边三角形，且\(DC=DE=CE=1\).

折叠后，长度保持不变，且\(A\)，\(B\)重合于\(P\)，从而\(PD=PE=PC=1\). 因此\(P-DCE\)的六条棱均为\(1\)，是正四面体. 正四面体棱长为\(1\)时，外接球半径为\(R=\dfrac{\sqrt{6}}{4}\)，故外接球体积为
\[
\frac43\pi R^3
=\frac43\pi\left(\frac{\sqrt6}{4}\right)^3
=\frac{\sqrt6\pi}{8}
\]
故选 C.
\end{solution}

\section{填空题}

\begin{problem}
若\(\displaystyle\lim_{n\to\infty}\dfrac{1}{\sqrt{n}\left(\sqrt{n+a}-\sqrt n\right)}=1\)，则常数\(a=\)\fillinblank{}.
\end{problem}

\begin{answer}
\(2\)
\end{answer}

\begin{solution}
由于分母含有\(a\)，必有\(a\ne0\). 有理化分母，得
\[
\frac{1}{\sqrt n\left(\sqrt{n+a}-\sqrt n\right)}
=\frac{\sqrt{n+a}+\sqrt n}{a\sqrt n}
=\frac{\sqrt{1+\frac an}+1}{a}
\]
令\(n\to\infty\)，极限为\(\dfrac2a\). 由题设\(\dfrac2a=1\)，所以\(a=2\).
\end{solution}

\begin{problem}
已知抛物线\(y^2=4x\)，过点\(P(4,0)\)的直线与抛物线相交于\(A(x_1,y_1)\)，\(B(x_2,y_2)\)两点，则\(y_1^2+y_2^2\)的最小值是\fillinblank{}.
\end{problem}

\begin{answer}
\(32\)
\end{answer}

\begin{solution}
直线\(y=0\)只与抛物线交于顶点，不能给出两个交点，故所考虑的直线均可写成\(x=4+ky\). 代入\(y^2=4x\)，得
\(y^2-4ky-16=0\).
其两个根为\(y_1\)，\(y_2\)，所以
\(y_1+y_2=4k\)，\(y_1y_2=-16\).
因此
\[
y_1^2+y_2^2=(y_1+y_2)^2-2y_1y_2
=16k^2+32\geqslant32
\]
当\(k=0\)时，直线为\(x=4\)，与抛物线交于\((4,4)\)，\((4,-4)\)，等号成立，故最小值为\(32\).
\end{solution}

\begin{problem}
如图，已知在正三棱柱\(ABC-A_1B_1C_1\)的所有棱长都相等，\(D\)是\(A_1C_1\)的中点，则直线\(AD\)与平面\(B_1DC\)所成角的正弦值为\fillinblank{}.

\bitmapfigure[max width=5.8cm]{img/2006/shandong_science/q15_fig1.png}
\end{problem}

\begin{answer}
\(\dfrac45\)
\end{answer}

\begin{solution}
设棱长为\(1\)，建立空间直角坐标系，取\(A=(0,0,0)\)，\(B=(1,0,0)\)，\(C=\left(\dfrac12,\dfrac{\sqrt3}{2},0\right)\)，\(A_1=(0,0,1)\)，\(B_1=(1,0,1)\)，\(C_1=\left(\dfrac12,\dfrac{\sqrt3}{2},1\right)\). 于是\(D=\left(\dfrac14,\dfrac{\sqrt3}{4},1\right)\)，\(\overrightarrow{AD}=\left(\dfrac14,\dfrac{\sqrt3}{4},1\right)\).

取平面\(B_1DC\)内的两个方向向量\(\overrightarrow{CB_1}=\left(\dfrac12,-\dfrac{\sqrt3}{2},1\right)\)，\(\overrightarrow{CD}=\left(-\dfrac14,-\dfrac{\sqrt3}{4},1\right)\). 它们的一个法向量可取\(\bs n=\left(\sqrt3,3,\sqrt3\right)\)，
且\(|\bs n|=\sqrt{15}\)，\(|\overrightarrow{AD}|=\dfrac{\sqrt5}{2}\)，\(\overrightarrow{AD}\cdot\bs n=2\sqrt3\). 设直线\(AD\)与平面\(B_1DC\)所成角为\(\theta\)，则
\[
\sin\theta=\frac{|\overrightarrow{AD}\cdot\bs n|}{|\overrightarrow{AD}|\,|\bs n|}=\frac{2\sqrt3}{\frac{\sqrt5}{2}\sqrt{15}}=\frac45
\]
\end{solution}

\begin{problem}
下列四个命题中，真命题的序号有\fillinblank{}.（写出所有真命题的序号）

\begin{circlelist}
\item 将函数\(y=|x+1|\)的图象按向量\(\bs v=(-1,0)\)平移，得到的图象对应的函数表达式为\(y=|x|\)；
\item 圆\(x^2+y^2+4x+2y+1=0\)与直线\(y=\dfrac12x\)相交，所得弦长为\(2\)；
\item 若\(\sin(\alpha+\beta)=\dfrac12\)，\(\sin(\alpha-\beta)=\dfrac13\)，则\(\tan\alpha\cot\beta=5\)；
\item 如图，已知正方体\(ABCD-A_1B_1C_1D_1\)，\(P\)为底面\(ABCD\)内一动点，\(P\)到平面\(AA_1D_1D\)的距离与到直线\(CC_1\)的距离相等，则\(P\)点的轨迹是抛物线的一部分.
\end{circlelist}

\bitmapfigure[max width=5.8cm]{img/2006/shandong_science/q16_fig1.png}
\end{problem}

\begin{answer}
\(\circled{3}\)、\(\circled{4}\)
\end{answer}

\begin{solution}
命题\(\circled{1}\)中，图象向量平移\((-1,0)\)后，点\((u,|u+1|)\)变为\((u-1,|u+1|)\)，令\(x=u-1\)，新函数为\(y=|x+2|\)，不是\(y=|x|\)，所以命题\(\circled{1}\)是假命题.

命题\(\circled{2}\)中，圆可化为\((x+2)^2+(y+1)^2=4\)，
圆心为\((-2,-1)\)，半径为\(2\). 圆心满足\(y=\dfrac12x\)，所以该直线过圆心，所得弦为直径，长度为\(4\)，不是\(2\)，故命题\(\circled{2}\)是假命题.

命题\(\circled{3}\)中，
\(2\sin\alpha\cos\beta=\sin(\alpha+\beta)+\sin(\alpha-\beta)=\dfrac56\)，
\(2\cos\alpha\sin\beta=\sin(\alpha+\beta)-\sin(\alpha-\beta)=\dfrac16\).
因此\(\tan\alpha\cot\beta=\dfrac{\sin\alpha\cos\beta}{\cos\alpha\sin\beta}=\dfrac{5/12}{1/12}=5\)，
故命题\(\circled{3}\)是真命题.

设正方体棱长为\(a\)，取底面坐标\(A=(0,0)\)，\(B=(a,0)\)，\(C=(a,a)\)，\(D=(0,a)\)，令\(P=(x,y)\). 点\(P\)到平面\(AA_1D_1D\)的距离为\(x\)，到直线\(CC_1\)的距离为\(\sqrt{(x-a)^2+(y-a)^2}\). 相等时
\[
x^2=(x-a)^2+(y-a)^2
\Longleftrightarrow
(y-a)^2=2ax-a^2
\]
这是抛物线方程的一部分，故命题\(\circled{4}\)是真命题. 所以真命题的序号为\(\circled{3}\)，\(\circled{4}\).
\end{solution}

\section{解答题}

\begin{problem}
已知函数\(f(x)=A\sin^2(\omega x+\varphi)\)（\(A>0\)，\(\omega>0\)，\(0<\varphi<\dfrac{\pi}{2}\)），且\(y=f(x)\)的最大值为\(2\)，其图象相邻两对称轴间的距离为\(2\)，并过点\((1,2)\).

\begin{enumerate}
\item 求\(\varphi\)；
\item 计算\(f(1)+f(2)+\cdots+f(2008)\).
\end{enumerate}
\end{problem}

\begin{solution}
\begin{enumerate}
\item 因为\(\sin^2\)的最大值为\(1\)，所以\(A=2\). 函数\(\sin^2 t\)的相邻对称轴对应的自变量之差为\(\dfrac{\pi}{2}\)，故\(\dfrac{\pi}{2\omega}=2\)，
从而\(\omega=\dfrac{\pi}{4}\). 又\(f(1)=2\)，所以
\(2\sin^2\left(\dfrac{\pi}{4}+\varphi\right)=2\)，
即\(\sin^2\left(\dfrac{\pi}{4}+\varphi\right)=1\). 由\(0<\varphi<\dfrac{\pi}{2}\)，有\(\dfrac{\pi}{4}+\varphi\in\left(\dfrac{\pi}{4},\dfrac{3\pi}{4}\right)\)，故
\(\varphi=\dfrac{\pi}{4}\).

\item 由上面的参数，\(f(k)=2\sin^2\left(\dfrac{(k+1)\pi}{4}\right)\). 所以\(f(1)=2\)，\(f(2)=1\)，\(f(3)=0\)，\(f(4)=1\)，且每四项循环一次. 因为\(2008=4\times502\)，故\(f(1)+f(2)+\cdots+f(2008)=502(2+1+0+1)=2008\).
\end{enumerate}
\end{solution}

\begin{problem}
设函数\(f(x)=ax-(a+1)\ln(x+1)\)，其中\(a\geqslant-1\)，求\(f(x)\)的单调区间.
\end{problem}

\begin{solution}
函数定义域为\((-1,+\infty)\)，且\(f'(x)=a-\dfrac{a+1}{x+1}=\dfrac{ax-1}{x+1}\).
当\(-1\leqslant a\leqslant0\)时，对任意\(x>-1\)，都有\(ax-1<0\)，而\(x+1>0\)，所以\(f'(x)<0\)，函数在\((-1,+\infty)\)上单调递减.

当\(a>0\)时，\(f'(x)=0\)的唯一解为\(x=\dfrac1a\)，并且在\(\left(-1,\dfrac1a\right)\)上\(f'(x)<0\)，在\(\left(\dfrac1a,+\infty\right)\)上\(f'(x)>0\).
因此当\(a>0\)时，函数在\(\left(-1,\dfrac1a\right)\)上单调递减，在\(\left(\dfrac1a,+\infty\right)\)上单调递增.
\end{solution}

\begin{problem}
如图，已知平面\(A_1B_1C_1\)平行于三棱锥\(V-ABC\)的底面\(ABC\)，等边\(\triangle AB_1C\)所在平面与底面\(ABC\)垂直，且\(\angle ACB=90^\circ\)，设\(AC=2a\)，\(BC=a\).

\begin{enumerate}
\item 求证直线\(B_1C_1\)是异面直线\(AB_1\)与\(A_1C_1\)的公垂线；
\item 求点\(A\)到平面\(VBC\)的距离；
\item 求二面角\(A-VB-C\)的大小.
\end{enumerate}

\bitmapfigure[max width=5.8cm]{img/2006/shandong_science/q19_fig1.png}
\end{problem}

\begin{solution}
建立直角坐标系，令\(C=(0,0,0)\)，\(A=(2a,0,0)\)，\(B=(0,a,0)\). 因为\(\triangle AB_1C\)为边长\(2a\)的等边三角形，且其所在平面垂直于底面，所以可取\(B_1=(a,0,\sqrt3a)\).

\begin{enumerate}
\item 由\(A_1B_1C_1\parallel ABC\)，得
\(B_1C_1\parallel BC\)，\(A_1C_1\parallel AC\).
又
\(\overrightarrow{AB_1}=(-a,0,\sqrt3a)\)，\(\overrightarrow{BC}=(0,-a,0)\)，
所以\(AB_1\perp BC\)，从而\(AB_1\perp B_1C_1\). 由于\(AC\perp BC\)，又有\(A_1C_1\perp B_1C_1\). 直线\(B_1C_1\)分别在\(B_1\)，\(C_1\)处与\(AB_1\)，\(A_1C_1\)相交，因此它是这两条异面直线的公垂线.

\item 因为\(B_1\)在棱\(VB\)上，平面\(VBC\)就是平面\(B_1BC\). 由\(C\)，\(B\)，\(B_1\)三点坐标，平面\(VBC\)的方程为
\(z=\sqrt3x\).
点\(A=(2a,0,0)\)到该平面的距离为
\(\dfrac{|\sqrt3\cdot2a-0|}{\sqrt{(\sqrt3)^2+(-1)^2}}=\sqrt3a\).

\item 平面\(AVB\)就是平面\(ABB_1\)，平面\(BVC\)就是平面\(BB_1C\). 取这两个平面的法向量
\[
\bs n_1=\overrightarrow{BA}\times\overrightarrow{BB_1}
=a^2(-\sqrt3,-2\sqrt3,-1)
\]
\[
\bs n_2=\overrightarrow{BC}\times\overrightarrow{BB_1}
=a^2(-\sqrt3,0,1)
\]
于是\(|\bs n_1|=4a^2\)，\(|\bs n_2|=2a^2\)，\(\bs n_1\cdot\bs n_2=2a^4\).
设二面角\(A-VB-C\)为\(\theta\)，则
\[
\cos\theta
=\frac{|\bs n_1\cdot\bs n_2|}{|\bs n_1|\,|\bs n_2|}
=\frac14
\]
故二面角的大小为\(\arccos\dfrac14\).
\end{enumerate}
\end{solution}

\begin{problem}
袋中装有标有数字\(1\)，\(2\)，\(3\)，\(4\)，\(5\)的小球各\(2\)个，从袋中任取\(3\)个小球，按\(3\)个小球上最大数字的\(9\)倍计分，每小球被取出的可能性都相等，用\(\xi\)表示取出的\(3\)个小球上的最大数字，求：

\begin{enumerate}
\item 取出的\(3\)个小球上的数字互不相同的概率；
\item 随机变量\(\xi\)的概率分布和数学期望；
\item 计分介于\(20\)分到\(40\)分之间的概率.
\end{enumerate}
\end{problem}

\begin{solution}
\begin{enumerate}
\item 总的取法数为\(\mathrm C_{10}^{3}=120\). 数字互不相同时，先从\(5\)个数字中选\(3\)个，再从每个数字的\(2\)个小球中各取\(1\)个，取法数为\(\mathrm C_{5}^{3}\cdot2^3=80\).
所以所求概率为\(\dfrac{80}{120}=\dfrac23\).

\item 最大数字为\(k\)时，先从标号不超过\(k\)的\(2k\)个球中取\(3\)个，再排除全部来自前\(k-1\)个数字的取法，取法数为\(\mathrm C_{2k}^{3}-\mathrm C_{2k-2}^{3}\). 因而随机变量\(\xi\)的概率分布为：\(P(\xi=1)=0\)，\(P(\xi=2)=\dfrac1{30}\)，\(P(\xi=3)=\dfrac2{15}\)，\(P(\xi=4)=\dfrac3{10}\)，\(P(\xi=5)=\dfrac8{15}\).
并且
\[
E(\xi)
=2\cdot\frac1{30}+3\cdot\frac2{15}
 +4\cdot\frac3{10}+5\cdot\frac8{15}
=\frac{13}{3}
\]

\item 得分为\(9\xi\). 由于\(20<9\xi<40\)时\(\xi=3\)，\(4\)，所求概率为\(P(\xi=3)+P(\xi=4)=\dfrac2{15}+\dfrac3{10}=\dfrac{13}{30}\).
\end{enumerate}
\end{solution}

\begin{problem}
双曲线\(C\)与椭圆\(\dfrac{x^2}{8}+\dfrac{y^2}{4}=1\)有相同的焦点，直线\(y=\sqrt3x\)为\(C\)的一条渐近线.

\begin{enumerate}
\item 求双曲线\(C\)的方程；
\item 过点\(P(0,4)\)的直线\(l\)，交双曲线\(C\)于\(A\)，\(B\)两点，交\(x\)轴于\(Q\)点（\(Q\)点与\(C\)的顶点不重合）. 当\(\overrightarrow{PQ}=\lambda_1\overrightarrow{QA}=\lambda_2\overrightarrow{QB}\)，且\(\lambda_1+\lambda_2=-\dfrac83\)时，求\(Q\)点的坐标.
\end{enumerate}
\end{problem}

\begin{solution}
\begin{enumerate}
\item 椭圆的焦半距满足
\(c^2=8-4=4\)，
所以双曲线的半焦距为\(c=2\). 设双曲线方程为\(\dfrac{x^2}{a^2}-\dfrac{y^2}{b^2}=1\)，则
\(a^2+b^2=c^2=4\). 渐近线斜率为\(\dfrac ba=\sqrt3\)，所以\(b^2=3a^2\). 因而\(4a^2=4\)，即\(a=1\)，\(b^2=3\). 双曲线方程为\(x^2-\dfrac{y^2}{3}=1\).

\item 设\(Q=(q,0)\). 因为直线过\(P=(0,4)\)，其上任一点可表示为
\[
(x,y)=Q+t\overrightarrow{PQ}=(q(1+t),-4t)
\]
其中\(t=0\)对应\(Q\)，\(t=-1\)对应\(P\). 设\(A\)，\(B\)对应的参数分别为\(t_1\)，\(t_2\). 由\(\overrightarrow{PQ}=\lambda_1\overrightarrow{QA}=\lambda_2\overrightarrow{QB}\)，得
\(t_1=\dfrac1{\lambda_1}\)，\(t_2=\dfrac1{\lambda_2}\).
将参数式代入双曲线方程，得\(q^2(1+t)^2-\dfrac{16t^2}{3}=1\)，即
\[
\left(q^2-\frac{16}{3}\right)t^2+2q^2t+q^2-1=0
\]
因\(Q\)点不与双曲线的顶点重合，\(q^2\ne1\)；又直线与双曲线有两个交点，故二次项系数不能为\(0\)，即\(q^2\ne\dfrac{16}{3}\)，可以使用韦达定理，
\(t_1+t_2=-\dfrac{2q^2}{q^2-\frac{16}{3}}\)，\(t_1t_2=\dfrac{q^2-1}{q^2-\frac{16}{3}}\).
因此\(\lambda_1+\lambda_2=\dfrac1{t_1}+\dfrac1{t_2}=\dfrac{t_1+t_2}{t_1t_2}=-\dfrac{2q^2}{q^2-1}\).
由题设
\[
-\frac{2q^2}{q^2-1}=-\frac83
\]
解得\(q^2=4\)，故\(q=\pm2\). 当\(q^2=4\)时，上述关于\(t\)的方程为\(-\dfrac43t^2+8t+3=0\)，其判别式为\(80>0\)，且两根乘积为\(-\dfrac94\ne0\)，所以这两个\(q\)值确实都对应直线与双曲线的两个不同交点，且不经过顶点. 因此\(Q=(2,0)\)或\(Q=(-2,0)\).
\end{enumerate}
\end{solution}

\begin{problem}
已知\(a_1=2\)，点\((a_n,a_{n+1})\)在函数\(f(x)=x^2+2x\)的图象上，其中\(n=1\)，\(2\)，\(3\)，\(\cdots\).

\begin{enumerate}
\item 证明数列\(\{\lg(1+a_n)\}\)是等比数列；
\item 设\(T_n=(1+a_1)(1+a_2)\cdots(1+a_n)\)，求\(T_n\)及数列\(\{a_n\}\)的通项；
\item 记\(b_n=\dfrac1{a_n}+\dfrac1{a_n+2}\)，求数列\(\{b_n\}\)的前\(n\)项和\(S_n\)，并证明\(S_n+\dfrac{2}{3T_n-1}=1\).
\end{enumerate}
\end{problem}

\begin{solution}
\begin{enumerate}
\item 由点\((a_n,a_{n+1})\)在函数图象上，\(a_{n+1}=a_n^2+2a_n\). 于是\(1+a_{n+1}=(1+a_n)^2\).
由\(a_1=2>0\)以及递推式\(a_{n+1}=a_n(a_n+2)\)可归纳证明\(a_n>0\)：若\(a_n>0\)，则\(a_{n+1}=a_n(a_n+2)>0\). 所以所有\(1+a_n>0\)，可以取常用对数，并由\(1+a_{n+1}=(1+a_n)^2\)得到\(\lg(1+a_{n+1})=2\lg(1+a_n)\).
故数列\(\{\lg(1+a_n)\}\)是首项\(\lg3\)、公比\(2\)的等比数列.

\item 由上一问，\(\lg(1+a_n)=2^{n-1}\lg3\)，所以\(1+a_n=3^{2^{n-1}}\)，\(a_n=3^{2^{n-1}}-1\).
进一步，
\[
T_n=\prod_{k=1}^{n}3^{2^{k-1}}
=3^{1+2+\cdots+2^{n-1}}
=3^{2^n-1}
\]

\item 由递推关系\(a_{n+1}=a_n(a_n+2)\)，有\(b_n=\dfrac1{a_n}+\dfrac1{a_n+2}=\dfrac{2(a_n+1)}{a_n(a_n+2)}=2\left(\dfrac1{a_n}-\dfrac1{a_{n+1}}\right)\). 所以\(S_n=\sum_{k=1}^{n}b_k=\dfrac2{a_1}-\dfrac2{a_{n+1}}=1-\dfrac2{a_{n+1}}\). 另一方面，\(3T_n-1=3^{2^n}-1=a_{n+1}\). 故\(S_n+\dfrac2{3T_n-1}=1-\dfrac2{a_{n+1}}+\dfrac2{a_{n+1}}=1\).
\end{enumerate}
\end{solution}
